% !TeX root = ../main.tex

\chapter{系统实现}
\label{cha:implementation}

本章详细阐述vlm-chat系统的工程实现，包括整体架构、六大核心模块以及关键工程亮点。

\section{系统架构}

vlm-chat采用\textbf{三层架构}设计，如图\ref{fig:architecture}所示。系统通过清晰的层次分离实现了关注点隔离，各层之间通过标准接口通信。

\begin{figure}[htbp]
  \centering
  \fbox{\parbox{0.88\textwidth}{
    \centering
    \vspace{5pt}
    \textbf{┌─────────────────────────────────────────┐}\\[2pt]
    \textbf{│\quad 表现层（Presentation Layer）\qquad\qquad\qquad\ \ │}\\[2pt]
    \textbf{├─────────────────────────────────────────┤}\\[2pt]
    \quad Gradio Web UI（Python）\quad | \quad React TypeScript SPA\\[2pt]
    \quad 图片上传 → 文本输入 → 流式聊天展示\\[2pt]
    \quad 会话列表管理 → Provider切换 → System Prompt编辑\\[3pt]
    \textbf{└─────────────────────────────────────────┘}\\[3pt]
    \qquad\qquad\qquad\qquad $\downarrow$ REST + SSE Streaming\\[3pt]
    \textbf{┌─────────────────────────────────────────┐}\\[2pt]
    \textbf{│\quad 业务逻辑层（Business Logic Layer）\qquad\qquad\ \ │}\\[2pt]
    \textbf{├─────────────────────────────────────────┤}\\[2pt]
    \quad api\_client.py（多Provider封装） → chat\_manager.py（会话管理）\\[2pt]
    \quad image\_processor.py（图片/PDF处理） → rag.py（知识库检索）\\[2pt]
    \quad tools.py（联网搜索） → embeddings.py（向量嵌入）\\[3pt]
    \textbf{└─────────────────────────────────────────┘}\\[3pt]
    \qquad\qquad\qquad\qquad $\downarrow$ OpenAI兼容API\\[3pt]
    \textbf{┌─────────────────────────────────────────┐}\\[2pt]
    \textbf{│\quad 模型服务层（Model Service Layer）\qquad\qquad\ \ │}\\[2pt]
    \textbf{├─────────────────────────────────────────┤}\\[2pt]
    \quad DashScope（Qwen3-VL）\ |\ OpenAI\ |\ OpenRouter\ |\ Ollama\\[2pt]
    \quad 云端API调用 · 无需本地GPU\\[3pt]
    \textbf{└─────────────────────────────────────────┘}
    \vspace{5pt}
  }}
  \caption{系统三层架构图}
  \label{fig:architecture}
\end{figure}

\subsection{各层职责}

\begin{itemize}
  \item \textbf{表现层}：提供Gradio（Python原生）和React TypeScript（Vite构建）双前端，均支持图片上传、流式输出和多轮对话。Gradio界面采用ChatGPT风格暗色主题双栏布局，React前端提供更灵活的组件化交互；
  \item \textbf{业务逻辑层}：Python后端核心，由6个模块和1个全局配置（config.py）组成，负责API调用封装、会话生命周期管理、图片验证与预处理、RAG知识库检索、工具调用编排等全部业务逻辑；
  \item \textbf{模型服务层}：通过OpenAI兼容接口统一接入多个模型Provider，支持DashScope（Qwen-VL系列）、OpenAI、OpenRouter、Ollama本地部署共4类后端，可在运行时热切换。
\end{itemize}

\section{核心模块设计}

系统业务逻辑层由6个核心模块构成，各模块职责明确、低耦合高内聚。

\subsection{多Provider API封装（api\_client.py）}

\texttt{api\_client.py}是系统与外部模型服务交互的统一入口，提供以下关键能力：

\begin{itemize}
  \item \textbf{多Provider支持}：通过PROVIDERS字典注册4类后端（DashScope、OpenAI、OpenRouter、Ollama），每种Provider预配置API地址、模型列表和默认参数。运行时可通过\texttt{switch\_model()}方法无感切换；
  \item \textbf{流式输出}：基于Python Generator模式实现SSE流式生成，\texttt{call\_api\_stream()}方法逐token产出响应，前端实时渲染；
  \item \textbf{消息构建}：\texttt{build\_messages()}方法根据是否有图片自动切换content格式——纯文本用字符串，带图片用\texttt{[\{type: "text", ...\}, \{type: "image\_url", ...\}]}列表，避免OpenAI API的400错误；
  \item \textbf{工具调用}：\texttt{call\_api\_with\_tools()}方法支持Function Calling，将DuckDuckGo搜索结果注入模型上下文进行二次生成；
  \item \textbf{容错机制}：自动重试3次（间隔2秒），网络波动或限流时自动恢复。
\end{itemize}

\subsection{会话管理（chat\_manager.py）}

\texttt{chat\_manager.py}负责会话全生命周期管理：

\begin{itemize}
  \item \textbf{会话CRUD}：创建、切换、删除、重命名会话，每个会话维护独立的消息历史和配置（Provider、模型、System Prompt）；
  \item \textbf{SQLite持久化}：会话元数据和消息历史写入本地SQLite数据库，支持应用重启后恢复全部对话上下文；
  \item \textbf{审计日志}：所有消息附带时间戳和模型标识，支持对话历史回溯；
  \item \textbf{上下文窗口管理}：自动截断过长历史，保留最近N轮对话以确保不超出模型token限制。
\end{itemize}

\subsection{图片与文档处理（image\_processor.py）}

\texttt{image\_processor.py}提供上传文件的验证与预处理：

\begin{itemize}
  \item \textbf{格式验证}：支持的图片格式包括JPEG、PNG、GIF、WebP、BMP，文档支持PDF。通过Pillow和PyPDF2进行格式检测；
  \item \textbf{大小限制}：图片最大10MB，PDF最大20页，超过限制时自动拒绝并提示用户；
  \item \textbf{Base64编码}：将图片编码为base64 data URI格式，直接嵌入API请求的content数组中；
  \item \textbf{临时文件管理}：上传文件存储至临时目录，按配置的保留时长（默认7天）自动清理。
\end{itemize}

\subsection{RAG知识库（rag.py）}

\texttt{rag.py}实现了基于本地向量检索的文档问答能力：

\begin{itemize}
  \item \textbf{文档入库}：上传PDF后，使用jieba分词将文本切分为语义块，调用\texttt{embeddings.py}生成哈希向量，存入SQLite向量表；
  \item \textbf{检索流程}：用户提问时，将问题向量化后在向量表中进行余弦相似度检索，返回Top-K（默认3）个最相关文档块；
  \item \textbf{上下文增强}：检索结果拼接至System Prompt中，使模型能基于文档内容回答，实现"上传PDF → 提问"的RAG闭环。
\end{itemize}

\subsection{工具调用（tools.py）}

\texttt{tools.py}实现了DuckDuckGo联网搜索的工具调用能力：

\begin{itemize}
  \item \textbf{搜索接口}：调用DuckDuckGo Instant Answer API，返回摘要和链接；
  \item \textbf{Function Calling编排}：模型判断需要联网时生成tool\_call，系统执行搜索后将结果作为新消息注入上下文，模型基于搜索结果进行二次生成；
  \item \textbf{用户可控}：UI中提供"启用联网搜索"开关，用户按需开启。
\end{itemize}

\subsection{向量嵌入（embeddings.py）}

\texttt{embeddings.py}提供轻量级的文本向量化能力：

\begin{itemize}
  \item \textbf{哈希嵌入}：基于字符n-gram和局部敏感哈希（LSH）生成固定维度向量，无需加载深度学习模型；
  \item \textbf{零依赖}：纯Python实现，不依赖transformers或sentence-transformers等重型库；
  \item \textbf{适用场景}：适合轻量级快速原型和中英文混合文本的关键词级检索。
\end{itemize}

\section{系统亮点}

vlm-chat相比基础VLM问答Demo，在工程层面实现了以下亮点功能。

\subsection{流式输出（Streaming）}

传统的"请求-等待-完整返回"模式在VLM场景下用户感知延迟高。vlm-chat采用\textbf{Server-Sent Events（SSE）}协议实现流式输出：API返回的每个token通过Generator逐块yield给前端，Gradio Chatbot组件实时增量渲染。用户看到文字逐字出现，大幅改善了交互体验。实测流式模式相比非流式模式的"首字延迟"缩短约60\%。

\subsection{多Provider热切换}

系统支持在DashScope（Qwen-VL系列）、OpenAI（GPT-4V等）、OpenRouter（聚合平台）和Ollama（本地部署）四类Provider之间运行时切换。切换时仅需更新API客户端实例的Provider标识和对应配置，会话上下文不丢失。该设计使用户可根据任务需求（中文OCR用DashScope、复杂推理用GPT-4V、隐私敏感用本地Ollama）灵活选择最优模型。

\subsection{RAG知识库检索}

对于长文档PDF，单次API调用的token限制难以覆盖全文。vlm-chat通过"PDF分块 → 向量嵌入 → 相似度检索 → 上下文注入"的RAG流程，使模型能基于检索到的相关文档片段进行问答。以20页PDF为例，RAG模式将有效上下文从约4K tokens提升至约12K tokens（3个相关块 + 原文上下文），回答质量显著优于直接截断的朴素方案。

\subsection{DuckDuckGo联网工具调用}

当用户问题涉及实时信息（如"今天的天气"、"最新的XX新闻"）时，模型知识截止日期可能成为限制。vlm-chat通过Function Calling机制集成DuckDuckGo搜索：模型自主判断是否需要联网，生成\texttt{web\_search}工具调用，系统执行搜索后将结果反馈模型进行二次推理。该流程实现了"静态模型知识 + 动态网络信息"的混合推理。

\subsection{限流与生产鉴权}

面向实际部署场景，系统内置了请求频率限制（REQUESTS\_PER\_MINUTE，默认20次/分钟）和Gradio鉴权（用户名/密码），防止API滥用和未授权访问。限流采用滑动窗口计数实现，溢出请求自动排队等待。

\subsection{双前端架构}

系统同时提供Gradio（Python原生，适合快速原型）和React TypeScript（Vite构建，适合生产部署）两套前端。两者共享同一后端API（FastAPI），React前端提供更灵活的组件化定制能力，包括主题切换、响应式布局等现代化特性。

\section{技术栈总览}

\begin{table}[htbp]
  \centering
  \caption{vlm-chat技术栈}
  \label{tab:tech_stack}
  \begin{tabular}{|l|l|l|}
    \hline
    \textbf{层次} & \textbf{技术选型} & \textbf{说明} \\
    \hline
    模型服务 & DashScope API (Qwen3-VL-Flash) & 默认Provider，中文优化 \\
    \hline
    后端框架 & Python 3.10 + FastAPI & REST API + SSE流式 \\
    \hline
    前端（Python） & Gradio 6.x & ChatGPT风格暗色主题 \\
    \hline
    前端（TS） & React 18 + Vite + TypeScript & 组件化生产级SPA \\
    \hline
    数据存储 & SQLite & 会话、消息、向量三表 \\
    \hline
    中文分词 & jieba & RAG分块 + 评测匹配 \\
    \hline
    图片处理 & Pillow + PyPDF2 & 格式验证与编码 \\
    \hline
    工具集成 & DuckDuckGo API & 联网搜索 \\
    \hline
    测试 & pytest (61用例) & Mock隔离外部依赖 \\
    \hline
    容器化 & Docker + docker-compose & 一键部署 \\
    \hline
  \end{tabular}
\end{table}
